\documentclass[11pt,a4paper]{article}
\usepackage{fontspec}
\usepackage{xeCJK}
\usepackage{graphicx}
\usepackage{xcolor}
\usepackage{hyperref}
\usepackage{titlesec}
\usepackage{enumitem}
\usepackage{fancyhdr}
\usepackage[a4paper,left=1.7cm,right=1.7cm,top=2.4cm,bottom=1.9cm,headheight=28pt,headsep=0.4cm]{geometry}
\setmainfont[Path=fonts/,Extension=.otf,BoldFont=*-Bold,ItalicFont=*,ItalicFeatures={FakeSlant=0.15},BoldItalicFont=*-Bold,BoldItalicFeatures={FakeSlant=0.15}]{NotoSerifSC}
\setCJKmainfont[Path=fonts/,Extension=.otf,BoldFont=*-Bold]{NotoSerifSC}
\hypersetup{hidelinks,pdftitle={陈宪章 - 学术简历},pdfauthor={Xian-Zhang Chen}}
\definecolor{primary}{RGB}{176,50,40}
\setlength{\parindent}{0pt}
\setlength{\parskip}{0.25em}
\linespread{1.08}
\setlength{\emergencystretch}{3em}
\setlist[itemize]{leftmargin=1.4em,itemsep=0.2em,topsep=0.3em}
\titleformat{\section}{\large\bfseries\color{primary}}{}{0em}{}[\titlerule]
\titlespacing*{\section}{0pt}{0.9em}{0.45em}
\titleformat{\subsection}{\normalsize\bfseries}{}{0em}{}
\titlespacing*{\subsection}{0pt}{0.65em}{0.3em}
\pagestyle{fancy}
\fancyhf{}
\renewcommand{\headrulewidth}{0.4pt}
\renewcommand{\footrulewidth}{0pt}
\fancyhead[L]{\includegraphics[height=0.7cm]{images/eit_logo.png}}
\fancyhead[R]{\small\color{primary}陈宪章\quad 学术简历}
\fancyfoot[L]{\footnotesize 更新：2026-10-26}
\fancyfoot[R]{\footnotesize\thepage}

\begin{document}
% Contact and employment: synchronize with data/profile.json.
{\LARGE\bfseries 陈宪章}\hfill Xian-Zhang Chen\par
Research Assistant Professor\par
Ningbo Institute of Digital Twin / 宁波东方理工大学\par
邮箱：\href{mailto:xzchen@idt.eitech.edu.cn}{xzchen@idt.eitech.edu.cn}\quad
\href{mailto:chenxzhlzu@gmail.com}{chenxzhlzu@gmail.com}\par
主页：\href{https://xianzhangchen.github.io/}{xianzhangchen.github.io}\quad
ORCID：\href{https://orcid.org/0000-0001-8821-7690}{0000-0001-8821-7690}\par
办公室：数理楼 306，宁波东方理工大学\par

\section{个人信息}
出生年月：1993.09\qquad 性别：男\qquad 政治面貌：中共党员

\section{工作经历}
\textbf{Ningbo Institute of Digital Twin / 宁波东方理工大学}\hfill 2025.12--至今\par
Research Assistant Professor；周通课题组。

\textbf{东方理工高等研究院 / 宁波东方理工大学}\hfill 2023.09--2025.11\par
博士后，凝聚态理论；合作导师：周通。

\section{教育背景}
\textbf{兰州大学，博士}\hfill 2016.09--2023.06\par
理论物理；导师：黄亮教授。

\textbf{斯特拉斯堡大学，联合培养博士}\hfill 2020.12--2022.11\par
IPCMS，凝聚态理论；导师：Rodolfo A. Jalabert 教授。

\textbf{兰州大学，本科}\hfill 2012--2016\par
理论物理。

\section{研究方向}
研究量子输运与拓扑量子器件，关注 Majorana 模、磁性织构与交错磁体。
结合解析模型与数值模拟，研究微观电子结构、拓扑性质和可测量输运信号之间的联系。
此前的研究包括石墨烯纳米结构中的扫描栅显微技术、量子疤痕及多体能谱统计。

\section{奖项荣誉}
\begin{itemize}
\item 入选欧洲博士学院（Collège doctoral Européen）国际博士生项目（PDI），2021 年。
\item 中国留学基金委奖学金（CSC），2020 年。
\item 第五届全国统计物理大会优秀海报奖，2019 年。
\item 硕士生国家奖学金，2016 年。
\end{itemize}

\newpage
\section{发表论文与预印本}
% BEGIN GENERATED PUBLICATIONS (scripts/sync_cv_data.py)
\subsection*{已发表论文}
\begin{itemize}[leftmargin=1.8em,label={},itemsep=0.55em,topsep=0.3em]
\item[\textbf{[9]}] Ziye Zhu, Richang Huang, \textbf{X.-Z. Chen}, Xunkai Duan, Jiayong Zhang, Igor Zutic, Tong Zhou, \textit{Altermagnetic Proximity Effect}, \href{https://doi.org/10.1103/kqy8-myz1}{Phys. Rev. Lett. 136, 186702 (2026)}.
\item[\textbf{[8]}] Zhou Cui, Ziye Zhu, Xunkai Duan, Bowen Hao, \textbf{X.-Z. Chen}, Jiayong Zhang, Tong Zhou, \textit{Symmetry-Driven Unconventional Magnetoelectric Coupling in Perovskite Altermagnets: From Bulk to the Two-Dimensional Limit}, \href{https://doi.org/10.1002/advs.202600004}{Advanced Science, e00004 (2026)}.
\item[\textbf{[7]}] Sihui Wu, Jiahui Qian, Xunkai Duan, \textbf{X.-Z. Chen}, Zhou Cui, Jiayong Zhang, Chunlan Ma, Tong Zhou, \textit{Electrically switchable quantum anomalous and spin Hall states in two-dimensional honeycomb lattices with non-Dirac bands}, \href{https://doi.org/10.1103/PhysRevB.111.235414}{Phys. Rev. B 111, 235414 (2025)}.
\item[\textbf{[6]}] \textbf{X.-Z. Chen}, Guillaume Weick, Dietmar Weinmann, Rodolfo A. Jalabert, \textit{Scanning gate microscopy in graphene nanostructures}, \href{https://doi.org/10.1103/PhysRevB.107.085420}{Phys. Rev. B 107, 085420 (2023)}.
\item[\textbf{[5]}] \textbf{X.-Z. Chen}, Zhen-Qi Chen, Liang Huang, Celso Grebogi, Ying-Cheng Lai, \textit{Many-body spectral statistics of relativistic quantum billiard systems}, \href{https://doi.org/10.1103/PhysRevResearch.5.013050}{Phys. Rev. Research 5, 013050 (2023)}.
\item[\textbf{[4]}] Xiao-Liang Li, \textbf{X.-Z. Chen}, Chen-Rong Liu, Liang Huang, \textit{Quantization condition of scarring states in complex soft-wall quantum billiards}, \href{https://doi.org/10.7498/aps.69.20200360}{Acta Phys. Sin. 69, 080506 (2020)}.
\item[\textbf{[3]}] Guan-Qun Zhang, \textbf{X.-Z. Chen}, Li Lin, Hai-Lin Peng, Zhong-Fan Liu, Liang Huang, Ning Kang, Hong-Qi Xu, \textit{Transport signatures of relativistic quantum scars in a graphene cavity}, \href{https://doi.org/10.1103/PhysRevB.101.085404}{Phys. Rev. B 101, 085404 (2020)}.
\item[\textbf{[2]}] Chen-Rong Liu, Pei Yu, \textbf{X.-Z. Chen}, Hong-Ya Xu, Liang Huang, Ying-Cheng Lai, \textit{Enhancing von Neumann entropy by chaos in spin-orbit entanglement}, \href{https://doi.org/10.1088/1674-1056/ab3dff}{Chin. Phys. B 28, 100501 (2019)}.
\item[\textbf{[1]}] Chen-Rong Liu, \textbf{X.-Z. Chen}, Hong-Ya Xu, Liang Huang, Ying-Cheng Lai, \textit{Effect of chaos on two-dimensional spin transport}, \href{https://doi.org/10.1103/PhysRevB.98.115305}{Phys. Rev. B 98, 115305 (2018)}.
\end{itemize}
\subsection*{预印本}
\begin{itemize}[leftmargin=1.8em,label={},itemsep=0.55em,topsep=0.3em]
\item[\textbf{[P3]}] \textbf{X.-Z. Chen}, Jiayong Zhang, Bowen Hao, Jiahui Qian, Ziye Zhu, Igor Zutic, Zhenyu Zhang, Tong Zhou, \textit{Altermagnets Enable Gate-Switchable Helical and Chiral Topological Transport with Spin-Valley-Momentum-Locked Dual Protection}, \href{https://arxiv.org/abs/2603.06487}{arXiv:2603.06487}.
\item[\textbf{[P2]}] Ziye Zhu, \textbf{X.-Z. Chen}, Xunkai Duan, Zhou Cui, Jiayong Zhang, Igor Zutic, Tong Zhou, \textit{Altermagnetoelectric Spin Field Effect Transistor}, \href{https://arxiv.org/abs/2512.02974}{arXiv:2512.02974}.
\item[\textbf{[P1]}] Richang Huang, Yongliang Hu, \textbf{X.-Z. Chen}, Peng Yu, Igor Zutic, Tong Zhou, \textit{Towards Scalable Braiding: Topological Superconductivity Unlocked under Nearly Arbitrary Magnetic Field Directions in Planar Josephson Junctions}, \href{https://arxiv.org/abs/2504.20031}{arXiv:2504.20031}.
\end{itemize}
% END GENERATED PUBLICATIONS

\newpage
\section{会议展示}
\subsection*{报告}
\begin{itemize}
\item Novel quantum spin and quantum anomalous Hall effects based on altermagnet.\par
CPS Meeting 2025, Harbin, China, 09/2025.
\item Many-body spectral statistics of relativistic quantum billiards systems.\par
DPG Meeting 2022, Regensburg, Germany, 09/2022.
\end{itemize}
\subsection*{海报}
\begin{itemize}
\item Transport signatures of scars in graphene dots at low magnetic field.\par
GDR 2426, Aussois, France, 12/2021.
\item Many-body spectral statistics of relativistic quantum billiard systems.\par
The 5th National Statistical Physics Conference (SPCSC), Hefei, China, 07/2019.
\item Transport signatures of relativistic quantum scars in a graphene cavity.\par
International Workshop on Experimental and Theoretical Developments on Complex Quantum Systems, Lanzhou, China, 07/2018.
\end{itemize}

\section{参与项目}
\begin{itemize}
\item Spectral statistics of a class of pseudo-integrable quantum billiards.\par
NSFC，项目编号 12175090，2021 年，经费 63 万元，参与。
\item Study of spin anomalous effects in Dirac billiards.\par
NSFC，项目编号 11775101，2017 年，经费 58 万元，参与。
\end{itemize}
\end{document}
